% Noether P08 Korean producer draft — UNCHECKED
% T04-U15; ED0002 lines6101--6110; source 586 LF B / C5E86BED66616D20E2C61BDA09965BDA2CCEA6034F0AAF0771438485D7548ADD
% Pointer v009 / ED0002: B06BE3530D9CF2E82B56FDBA7FE41D5D044DF2425DFA2A059D4939EAA2F7A6C2 / C9A125167ACB33D914EE4374B65AE7CDF0052F568371B8B77B720EA178ABF0E3
% Translation only; formulas are preserved for independent checking.

\emph{a) $\mathcal L$과 $\mathcal S$의 랭크를 비교하는 데에는 다음 보조정리 a)가 쓰인다. $f(x,y,\lambda_1x+\lambda_2y)=0$이면 반드시 $f(x,y,z)=0$이다.}

이는 세 개의 이진 변수열 사이에 성립하는 항등관계
\[
  (xy)z+(yz)x+(zx)y=0,
  \qquad \text{여기서 }(xy)=(x_1y_2-x_2y_1),\ldots,
\]
에서 곧바로 따른다. 이 관계로부터
\[
  f[x,y,(zy)x+(xz)y]=(xy)^p f(x,y,z)
\]
를 얻는다. 따라서 $f(x,y,\lambda_1x+\lambda_2y)=0$이 $\lambda$에 관한 항등식으로 성립하면, $\lambda_1=(zy)$, $\lambda_2=(xz)$를 대입하고 $(xy)\ne0$임을 이용하여 반드시 $f(x,y,z)=0$임을 얻는다.
